\section{Global governance: majority support for global democracy on world issues}

The preceding sections examined whether citizens accept specific forms of international redistribution. A distinct question is whether they also accept the supranational institutions needed to coordinate national contributions, limit free-riding, and implement common policies across countries.

% AF: useless
The evidence reviewed thus far indicates that citizens in surveyed countries may accept international redistribution through foreign aid, taxation, and other cross-border transfers. Public acceptance is only a first step, however, because these policies also require institutional arrangements capable of coordinating contributions and ensuring effective implementation. 
As \citet{kopczuk_limitations_2005} argue, decentralized national redistributive systems have limited effects on global inequality, while voluntary international transfers remain constrained by free-riding. Some degree of global coordination may therefore be necessary to reduce extreme poverty and address inequalities extending beyond national borders. This section examines whether citizens accept the forms of global governance required to overcome these coordination problems.

% AF: je vois pas l'intérêt du paragraphe suivant non plus, ça ne parle pas des attitudes
The need for coordination extends beyond international redistribution. Global inequality, climate change, pandemics, armed conflict, and migration create cross-border effects that national policies cannot address effectively in isolation. These challenges have generated an extensive literature on multilateral agreements, international organizations, and institutional reform. Global governance, however, need not entail a unitary world state. It can instead be understood as the rules, institutions, organizations, and decision-making procedures through which countries manage genuinely transnational problems. Under this conception, global governance does not replace domestic policymaking but addresses problems whose externalities and coordination requirements justify action above the national level \citep{ghassim_who_2024,dryzek_global_2011}.

% AF: pareil, ça parle pas d'attitudes
Public acceptance should therefore depend on the institutional design and scope of supranational authority. In particular, citizens may distinguish between an unspecified transfer of sovereignty and an institution that is democratically organized and restricted to problems requiring international coordination. \citet{dryzek_global_2011} provide a deliberative framework for evaluating such legitimacy, emphasizing the participation of affected actors, deliberation across competing perspectives, and effective transmission between civil society and authoritative decision-making. Global institutions must consequently be assessed not only by the geographical level at which they operate, but also by the democratic procedures through which authority is exercised. Their establishment ultimately depends on whether citizens accept both their mandate and their institutional design.

\citet{ghassim_who_2024} examine attitudes toward world government using survey data from approximately 42,000 respondents across 17 countries collected between 2017 and 2021. % AF: say which countries (perhaps in footnote)
Their experimental design varies two dimensions of legitimacy. Institutional legitimacy concerns whether the proposed global authority is explicitly democratic, whereas functional legitimacy concerns whether its mandate is confined to problems requiring global coordination. The results show that both dimensions substantially shape % AF: increase?
public acceptance.

With equal country weights, 52\% of respondents reject an unspecified world government, whereas population weighting yields 58\% support. When the proposed government is explicitly democratic, support rises to 67\% with equal country weights and 71\% with population weights. Restricting its mandate to global issues produces corresponding support rates of 66\% and 69\%, while a government focused specifically on COVID-19 receives 64\% and 69\% support. Combining institutional and functional legitimacy produces still higher acceptance: a democratic world government focused on global issues receives 69\% support with equal country weights and 72\% with population weights, while a democratic government focused on COVID-19 receives 72\% and 74\%, respectively \citep{ghassim_who_2024}.

The democratic global-issues proposal attracts majority support in every surveyed country except the United States. Support ranges from 75\% to 82\% in Egypt, India, Kenya, Indonesia, South Korea, Colombia, and Hungary; apart from the United States, the lowest levels are recorded in Russia, at 56\%, and Argentina, at 58\%. Citizens therefore respond not only to the transfer of authority beyond the nation-state but also to whether that authority is democratically constituted and functionally confined to genuinely global problems \citep{ghassim_who_2024}.

% AF: tu peux retirer ce paragraphe
The evidence thus traces a progressive expansion in the domain of acceptable redistribution. Universalistic values make concern for distant populations possible; attitudes toward foreign aid and international taxation show that this concern can extend to concrete transfers; and the findings of \citet{ghassim_who_2024} indicate that it may also extend to the institutions needed to govern transnational problems. The next question is whether this acceptance persists when respondents evaluate specific institutional reforms rather than an abstract model of global authority.

\citet{fabre_public_2026} examines whether acceptance of global governance extends beyond an abstract world government to concrete forms of international cooperation and institutional reform. When asked whether governments should cooperate to promote convergence in GDP per capita across countries by the end of the century, 61\% of respondents answer positively and 26\% negatively. Excluding non-responses, the overall acceptance rate reaches 70\%, although it is lower in the United States, at 56\%. Reform of the United Nations Security Council % AF: détailler quelle réforme
likewise receives 69\% acceptance, indicating that citizens may support not only internationally redistributive outcomes but also changes to the institutional architecture through which global decisions are made.

% AF: déjà dit plus haut, enlever
This acceptance may also have electoral implications. If the party to which they feel closest joined a global movement promoting climate action, millionaire taxation, and poverty reduction in low-income countries, 36\% of respondents report that they would become more likely to vote for it, compared with 17\% who would become less likely. Among non-voters who identify with a left-wing party, the corresponding figures are 46\% and 10\%. Although these responses concern hypothetical voting intentions, they indicate that global cooperation may have mobilizing potential even when it remains weakly salient in ordinary political competition \citep{fabre_public_2026}.

% AF: déjà dit plus haut, enlever
Support also persists when cooperation is embedded in a costly policy package. Overall, 68\% prefer a scenario combining global cooperation, millionaire taxation, warming limited to +2\textdegree{}C, vehicle electrification, and doubled prices for heating fuel, air travel, and beef to a status quo leading to +3\textdegree{}C warming by 2100 \citep{fabre_public_2026}. This finding demonstrates substantial acceptance of a coherent but costly package addressing climate change and global inequality, although it does not identify which component drives that support.

% AF: je reformulerais en mode: Une majorité est prête pour un gouvernement mondial, dès lors que celui-ci serait démocratique et limiterait sa compétence aux enjeux mondiaux, conformément au principe de subsidiarité. 
Taken together, these results complement the findings of \citet{ghassim_who_2024}. Public acceptance is greater when global governance is associated with clearly defined objectives, recognizable institutional reforms, and democratic procedures. The evidence does not establish unconditional demand for a centralized world government; rather, it indicates conditional acceptance of democratic, functionally limited, and policy-oriented forms of global governance \citep{fabre_public_2026,ghassim_who_2024}.

% AF: pas besoin de faire trop de paragraphes de liaison comme ça
Acceptance of global institutions in principle does not, however, establish that citizens will bear the material costs generated by the policies those institutions adopt. Climate policy provides a particularly demanding test because its private costs, distributional consequences, and dependence on participation by other countries can be made explicit.



